\def\writeANS{}
\begin{loigiaibth}{24}
  - Enthalpy tạo thành của một chất là nhiệt kèm theo phản ứng tạo thành $1$ mol chất đó từ các đơn chất bền nhất. \\ - Biến thiên Enthalpy của phản ứng là nhiệt lượng tỏa ra hay thu vào của một phản ứng hóa học trong quá trình đẳng áp (áp suất không đổi) \par \begin {tabular}{|l|l|} \hline \textbf {Enthalpy tạo thành của một chất} & \textbf {Biến thiên enthalpy của phản ứng} \\ \hline - Chất tham gia phải là đơn chất bền nhất. & - Chất tham gia ở dạng đơn chất hay hợp chất đều được. \\ \hline - Sản phẩm chỉ có $1$ chất duy nhất. & - Sản phẩm có thể là $1$ hay nhiều chất. \\ \hline - VD: & - VD: \\ $C$(graphite, s) + $O_2(g) \longrightarrow CO_2(g)$ & $C_2H_4(g) + H_2(g) \longrightarrow C_2H_6(g)$ \\ $\Delta _r H_{298}^0 = -393{,}5$ kJ / mol & $\Delta _r H_{298}^0 = -137{,}0$ kJ \\ \hline \end {tabular}
\end{loigiaibth}
\def\writeANS{}
\begin{loigiaibth}{25}
  \begin {enumerate} \item Nung $NH_4Cl(s)$ tạo ra $HCl(g)$ và $NH_3(g)$ là phản ứng thu nhiệt. \\ Giải thích: Do cần cung cấp nhiệt trong suốt quá trình phản ứng. \item Cồn cháy trong không khí là phản ứng tỏa nhiệt. \\ Giải thích: Do phản ứng đốt cháy cồn chỉ cần cung cấp nhiệt vào thời điểm ban đầu và có tỏa nhiệt trong quá trình phản ứng \item Phản ứng thủy phân collagen thành gelain (là một loại protein dễ tiêu hóa) diễn ra khi hầm xương động vật là phản ứng thu nhiệt. \\ Giải thích: Do cần cung cấp nhiệt trong suốt quá trình phản ứng (hầm). \end {enumerate}
\end{loigiaibth}
\def\writeANS{}
\begin{loigiaibth}{26}
  \centering \begin {tabular}{|c|l|c|c|} \hline STT & Quá trình & Tỏa nhiệt & Thu nhiệt \\ \hline 1 & Hóa hơi $X(l) \longrightarrow X(g)$ & & X \\ \hline 2 & Ngưng tụ $X(g) \longrightarrow X(l)$ & X & \\ \hline 3 & Thăng hoa $X(s) \longrightarrow X(g)$ & & X \\ \hline 4 & Nóng chảy $X(s) \longrightarrow X(l)$ & & X \\ \hline 5 & Đông đặc $X(l) \longrightarrow X(s)$ & X & \\ \hline \end {tabular}
\end{loigiaibth}
\def\writeANS{}
\begin{loigiaibth}{27}
  \begin {enumerate} \item Nước hóa rắn là quá trình tỏa nhiệt vì nước từ lỏng về rắn giảm nhiệt độ $\Rightarrow $ giải phóng nhiệt ra môi trường. \item Sự tiêu hóa thức ăn là là quá trình thu nhiệt vì cần cung cấp năng lượng để tiêu hóa thức ăn. \item Quá trình chạy của con người là quá trình tỏa nhiệt vì khi con người vận động sẽ xảy ra các phản ứng oxi hóa các chất dinh dưỡng như glucozơ sẽ tỏa nhiệt. \item Khí $CH_4$ đốt ở trong lò là quá trình tỏa nhiệt vì phản ứng chỉ cần đốt nóng lúc đầu, sau đó thì tự cháy (phát sáng và tỏa nhiệt). \item Hòa tan $KBr$ vào nước làm cho nước trở nên lạnh là quá trình thu nhiệt do làm nước lạnh đi. \item Thêm sulfuric acid đặc vào nước, nước nóng lên là quá trình tỏa nhiệt do làm nước nóng lên. \end {enumerate}
\end{loigiaibth}
\def\writeANS{}
\begin{loigiaibth}{28}
  Sơ đồ (1) là quá trình tỏa nhiệt, do nhiệt độ phản ứng tăng so với nhiệt độ ban đầu (nhiệt độ phòng) \\ Sơ đồ (2) là quá trình thu nhiệt, do nhiệt độ phản ứng giảm so với nhiệt độ ban đầu.
\end{loigiaibth}
\def\writeANS{}
\begin{loigiaibth}{29}
  Phản ứng (2) có $\Delta _r H_{298}^0$ âm hơn $\Rightarrow $ tỏa ra lượng nhiệt lớn hơn nên xảy ra thuận lợi hơn
\end{loigiaibth}
\def\writeANS{}
\begin{loigiaibth}{30}
  Đây là phản ứng thu nhiệt vì $\Delta _r H_{298}^0 > 0$ \\ Trong các sản phẩm tự nhiên, baking soda ($NaHCO_3$) giúp làm sạch, khử mùi, làm mềm mảng bám khi giấm ($CH_3COOH$) giúp loại bỏ mùi hôi và một số vết bẩn cứng đầu khác. \\ Ngoài tác dụng tẩy rửa của phản ứng giữa baking soda và giấm những ứng dụng khác của phản ứng trên là: trắng quần áo, thông bồn cầu, vệ sinh máy giặt, khử mùi...
\end{loigiaibth}
\def\writeANS{}
\begin{loigiaibth}{31}
  Khi phản ứng $2NH_4NO_3(s) + Ba(OH)_2.8H_2O(s) \longrightarrow 2NH_3(aq) + Ba(NO_3)_2(aq) + 10H_2O(l)$ diễn ra ở nhiệt độ phòng thì nhiệt độ của hỗn hợp bị giảm đi do phản ứng có $\Delta _r H_{298}^0 > 0$ là phản ứng thu nhiệt (hay hấp thụ năng lượng dưới dạng nhiệt) ở môi trường xung quanh làm nhiệt độ của hỗn hợp giảm (lạnh đi).
\end{loigiaibth}
\def\writeANS{}
\begin{loigiaibth}{32}
  Nhiệt tỏa ra khi hình thành $1$ mol $Na_2O(s)$ ở điều kiện chuẩn từ phản ứng giữa $Na(s)$ và $O_3$ $(g)$ không được coi là nhiệt tạo thành chuẩn của $Na_2O(s)$ vì $O_3$ $(g)$ không phải dạng bền nhất của đơn chất oxygen.
\end{loigiaibth}
\def\writeANS{}
\begin{loigiaibth}{33}
  \begin {enumerate} \item Do $\Delta _r H_{298}^0 < 0$ nên nước $(H_2O)$ có năng lượng thấp hơn hỗn hợp oxygen và hydrogen. \item Sơ đồ biến thiên năng lượng \begin {center} \begin {tikzpicture}[declare function={r=3;}, line join=round, line cap=round, >=stealth, scale=1] \draw [\mycolor !80!black,>=stealth,->,text=\maunhan !80!black] (0,0) -- (0,5.25) node [pos=0.9, font=\scriptsize \bfseries \sffamily ,anchor=east]{Enthalpy, kJ mol$^{-1}$}; \draw [\mycolor !80!black,>=stealth,->,text=\maunhan !80!black] (0,0) -- (6.25,0) ; \draw [dashed,\mycolor !80!black,text=\maunhan !80!black] (0,4)node[left,font=\scriptsize \bfseries \sffamily ]{$\Delta _{f}\text {H}^\circ _{298} \text {(cd)}$}--(3,4); \draw [dashed,\mycolor !80!black,text=\maunhan !80!black] (0,1)node[left,font=\scriptsize \bfseries \sffamily ]{$\Delta _{f}\text {H}^\circ _{298} \text {(sp)}$}--(6,1); \draw [line width=1.2pt,\mycolor !80!black,text=\maunhan !80!black] (0.5,4)--(3,4) node [pos=0.5, font=\scriptsize \bfseries \sffamily ,anchor=south]{$2\text {H}_2\text {(g) + O}_2\text {(g)}$} --(3,1) (3,4) edge[-stealth] (3,1) (3,1)--(6,1) node [pos=0.5, font=\scriptsize \bfseries \sffamily ,anchor=north]{$2\text {H}_2\text {O(l)}$}; \node at (4.5,2.5) [font=\scriptsize \bfseries \sffamily ] {$\Delta _{r}\text {H}^\circ _{298} = -484 \text {kJ}$}; \end {tikzpicture} \end {center} \end {enumerate}
\end{loigiaibth}
\def\writeANS{}
\begin{loigiaibth}{34}
  \begin {center} \begin {tikzpicture}[declare function={r=3;}, line join=round, line cap=round, >=stealth, scale=1] \draw [\mycolor !80!black,>=stealth,->,text=\maunhan !80!black] (0,0) -- (0,5.25) node [pos=0.9, font=\scriptsize \bfseries \sffamily ,anchor=east]{Năng lượng}; \draw [\mycolor !80!black,>=stealth,->,text=\maunhan !80!black] (0,0) -- (6.25,0) node [pos=1, font=\scriptsize \bfseries \sffamily ,anchor=north east]{Tiến trình phản ứng} ; \draw [dashed,\mycolor !80!black,text=\maunhan !80!black] (0,1)node[left,font=\scriptsize \bfseries \sffamily ]{$\Delta _{r}\text {H}^\circ _{298} \text {(cd)}$}--(3,1); \draw [dashed,\mycolor !80!black,text=\maunhan !80!black] (0,4)node[left,font=\scriptsize \bfseries \sffamily ]{$\Delta _{r}\text {H}^\circ _{298} \text {(sp)}$}--(6,4); \draw [line width=1.2pt,\mycolor !80!black,text=\maunhan !80!black] (0.5,1)--(3,1) node [pos=0.5, font=\scriptsize \bfseries \sffamily ,anchor=north]{$\text {CaCO}_3\text {(s)}$}--(3,4) (3,1) edge[-stealth] (3,4) (3,4)--(6,4) node [pos=0.5, font=\scriptsize \bfseries \sffamily ,anchor=south]{$\text {CaO(s) + CO}_2\text {(g)}$}; \node at (4.5,2.5) [font=\scriptsize \bfseries \sffamily ] {$\Delta _{r}\text {H}^\circ _{298} = +178{,}49 \text {kJ}$}; \end {tikzpicture} \hspace {0.5cm} \begin {tikzpicture}[declare function={r=3;}, line join=round, line cap=round, >=stealth, scale=1] \draw [\mycolor !80!black,>=stealth,->,text=\maunhan !80!black] (0,0) -- (0,5.25) node [pos=0.9, font=\scriptsize \bfseries \sffamily ,anchor=east]{Năng lượng}; \draw [\mycolor !80!black,>=stealth,->,text=\maunhan !80!black] (0,0) -- (6.25,0) node [pos=1, font=\scriptsize \bfseries \sffamily ,anchor=north east]{Tiến trình phản ứng} ; \draw [dashed,\mycolor !80!black,text=\maunhan !80!black] (0,4)node[left,font=\scriptsize \bfseries \sffamily ]{$\Delta _{r}\text {H}^\circ _{298} \text {(cd)}$}--(3,4); \draw [dashed,\mycolor !80!black,text=\maunhan !80!black] (0,1)node[left,font=\scriptsize \bfseries \sffamily ]{$\Delta _{r}\text {H}^\circ _{298} \text {(sp)}$}--(6,1); \draw [line width=1.2pt,\mycolor !80!black,text=\maunhan !80!black] (0.5,4)--(3,4) node [pos=0.5, font=\scriptsize \bfseries \sffamily ,anchor=south]{$\text {NaOH(aq) + HCl(aq)}$} --(3,1) (3,4) edge[-stealth] (3,1) (3,1)--(6,1) node [pos=0.5, font=\scriptsize \bfseries \sffamily ,anchor=north]{$\text {NaCl(aq) + H}_2\text {O(l)}$}; \node at (4.5,2.5) [font=\scriptsize \bfseries \sffamily ] {$\Delta _{r}\text {H}^\circ _{298} = -57{,}3 \text {kJ}$}; \end {tikzpicture} \end {center}
\end{loigiaibth}
\def\writeANS{}
\begin{loigiaibth}{35}
  (a) Số mol $\mathrm {SO}_2$: $1{,}125$ mol. \\ Lượng nhiệt giải phóng: $-98{,}5 \times 1{,}125 = -110{,}81$ kJ \\ (b) Giá trị $\Delta _r \mathrm {H}_{298}^0$ của phản ứng: $\mathrm {SO}_3\mathrm {(g)} \longrightarrow \mathrm {SO}_2\mathrm {(g)} + \frac {1}{2}\mathrm {O}_2\mathrm {(g)}$ là $98{,}5$ kJ
\end{loigiaibth}
\def\writeANS{}
\begin{loigiaibth}{36}
  Ta có: (2) = 3.(1) $\Rightarrow \Delta _r H_{298}^\circ (2) = 3.(-32{,}4) = -97{,}2 \mathrm {kJ}$. \\ (3) = $\dfrac {(1)}{2} \Rightarrow \Delta _r H_{298}^\circ (3) = \dfrac {\Delta _r H_{298}^\circ (1)}{2} = \dfrac {-32{,}4}{2} = -16{,}2 \mathrm {kJ}$. \\ (4) = $\dfrac {-(1)}{4} \Rightarrow \Delta _r H_{298}^\circ (4) = \dfrac {-\Delta _r H_{298}^\circ (1)}{4} = \dfrac {32{,}4}{4} = 8{,}1\mathrm {kJ}$.
\end{loigiaibth}
\def\writeANS{}
\begin{loigiaibth}{37}
  Theo đề bài đốt cháy hoàn toàn $1$ gam $C_2H_2(g)$ ở điều kiện chuẩn, thu được $CO_2(g)$ và $H_2O(l)$, giải phóng $49{,}98$ kJ. Mà $1$ mol $C_2H_2$ có khối lượng $26$ gam \\ $\rightarrow $ nhiệt lượng giải phóng khi đốt cháy $1$ mol $C_2H_2 = 26.49{,}98 = 1299{,}48$ kJ \\ $\rightarrow \Delta _r H_{298}^0 = -1299{,}48$ kJ (do đây là phản ứng tỏa nhiệt nên giá trị enthalpy mang giá trị âm)
\end{loigiaibth}
\def\writeANS{}
\begin{loigiaibth}{38}
  $2H_2(g) + O_2(g) \rightarrow 2H_2O(l) \quad \Delta _r H_{298}^\circ = -571{,}6$ kJ \\ $Cu(OH)_2(s) \xrightarrow {t^\circ } CuO(s) + H_2O(l) \quad \Delta _r H_{298}^\circ = +9{,}0$ kJ \\ $Cu(graphite) + O_2(g) \xrightarrow {t^\circ } CO_2(g) \quad \Delta _r H_{298}^\circ = -393{,}5$ kJ
\end{loigiaibth}
\def\writeANS{}
\begin{loigiaibth}{39}
  Số mol $N_2 = 7 : 28 = 0{,}25$ mol \\ Để tạo $1$ mol $NH_3$ cần $0{,}5$ mol $N_2 \Rightarrow \Delta _r H_{298}^\circ = -22{,}95 \times 2 = -45{,}9$ kJ \\ Phương trình nhiệt hóa học: $N_2(g) + 3H_2(g) \rightleftharpoons 2NH_3(g) \quad \Delta _r H_{298}^\circ = -91{,}8$ kJ
\end{loigiaibth}
\def\writeANS{}
\begin{loigiaibth}{40}
  Ta có: (3) = (1) - (2) $\Rightarrow \Delta _r H_{298}^\circ (3) = \Delta _r H_{298}^\circ (1) - \Delta _r H_{298}^\circ (2) = -1411 - (-1367) = -44$ kJ.
\end{loigiaibth}
\def\writeANS{}
\begin{loigiaibth}{41}
  Ta có: (4) = -(1) + 2.(2) + 3.(3) $\Rightarrow \Delta _r H_{298}^\circ (4) = -(-84{,}7) + 2.(-393{,}5) + 3.(-285{,}6) = -1559{,}1$ kJ.
\end{loigiaibth}
\def\writeANS{}
\begin{loigiaibth}{42}
  \begin {enumerate} \item [(a)] $2Al(s) + 3Cl_2(g) \rightarrow 2AlCl_3(s) \quad \Delta _r H_{298}^\circ = -1390{,}81$ kJ \\ Đây là phản ứng oxi hóa- khử vì có sự thay đổi số oxi hóa của các nguyên tử trong phản ứng. \item [(b)] $\Delta _r H_{298}^\circ = -1390{,}81$ kJ $ < 0$ phản ứng tỏa nhiệt \item [(c)] $n_{AlCl_3} = \dfrac {10}{133{,}5} = 0{,}075$ mol \\ Cứ $2$ mol $Al$ phản ứng tạo thành $2$ mol $AlCl_3$ thì giải phóng nhiệt lượng là $1390{,}81$ kJ \\ $0{,}075$ mol ----------------------------------- $x = \dfrac {0{,}075.1390{,}81}{2} = 52{,}2$kJ. \item [(d)] Cứ $2$ mol $Al$ phản ứng thì giải phóng nhiệt lượng là $1390{,}81$ kJ \\ $x = \dfrac {2.1}{1390{,}81} = 1{,}438.10^{-3}$ mol. $\leftarrow \quad 1$ kJ \\ $\Rightarrow m_{Al} = 0{,}0388$ gam \end {enumerate}
\end{loigiaibth}
\def\writeANS{}
\begin{loigiaibth}{43}
  \begin {enumerate} \item [(a)] Phản ứng đó tỏa nhiệt vì có biến thiên enthalpy âm \item [(b)] Phản ứng thủy phân đường sucrose môi trường acid và đun nóng: \\ $C_{12}H_{22}O_{11} + H_2O \rightarrow C_6H_{12}O_6 \text {(glucose)} + C_6H_{12}O_6 \text {(fructose)}$ \\ Phản ứng trong sơ đồ là phản ứng oxi hóa- khử, oxygen là chất oxi hóa, đường glucose và fructose là chất khử $C_6H_{12}O_6(s) + 6 \ O_2(g) \rightarrow 6 \ CO_2(g) + 6 \ H_2O(l)$ \item [(c)] Phản ứng đốt cháy đường sucrose: $C_{12}H_{22}O_{11}(s) + 12O_2(g) \rightarrow 12 \ CO_2(g) + 11 \ H_2O(l)$ \\ Biến thiên enthalpy chuẩn của phản ứng là -$5645$kJ. \item [(d)] Biến thiên enthalpy của quá trình $= \dfrac {5}{342}.(-5645) = -82{,}5$kJ. \item [(e)] Cơ thể cần năng lượng để hoạt động nên phải có chế độ dinh dưỡng đầy đủ. Luyện tập thể dục, thể thao hợp lí giúp đốt cháy năng lượng dư thừa trong cơ thể. \end {enumerate}
\end{loigiaibth}
\def\writeANS{}
\begin{loigiaibth}{44}
  - Lượng nhiệt cần để đưa $1$ gam ethanol từ $20^\circ C$ đến $78{,}29^\circ C$ là $Q = 1{,}44(78{,}29 - 20) = 83{,}9376$ $J$. \par - Lượng nhiệt cần để hóa hóa hơi $1$ gam ethanol là $855$ $J$. \par $\Rightarrow $ Lượng nhiệt cần để đưa $1$ $kg$ ethanol từ $20^\circ C$ đến nhiệt độ sôi và hóa hơi hoàn toàn là $(83{,}9376+855).10^3 = 9{,}39.10^5$ $J$
\end{loigiaibth}
\def\writeANS{}
\begin{loigiaibth}{45}
  Ta có phản ứng nung vôi: $CaCO_3(s) \rightarrow CaO(s) + CO_2(s) \,\,\, \Delta _r H_{298}^0 = 178{,}29kJ$ \par Tức là để tạo $1$ mol $CaO$ cần cung cấp $178{,}29$ $kJ$ $\rightarrow $ phản ứng đốt cháy $CH_4$ cần tỏa ra $178{,}29$ $kJ$ \par Mà ta có phản ứng đốt cháy $CH_4$ : $CH_4(g) + 2O_2(g) \rightarrow CO_2(g) + 2H_2O(l) \,\,\, \Delta _r H_{298}^0 = -890{,}36kJ$ \par Đốt cháy $1$ mol $CH_4$ tỏa ra $890{,}36$ $kJ$ \par Đốt cháy $x$ mol $CH_4$ tỏa ra $178{,}29$ $kJ$ \par $\rightarrow x = \dfrac {178{,}29.1}{890{,}36} \approx 0{,}2 \, mol \rightarrow $ khối lượng $CH_4$ cần dùng để đốt cháy là: $0{,}2.16 = 3{,}2$ gam.
\end{loigiaibth}
\def\writeANS{}
\begin{loigiaibth}{46}
  Trong $1800$ g than đá chứa $1620$ gam ($135$ mol) carbon và $21{,}6$ g ($0{,}675$ mol) sulfur. \par Nhiệt lượng giải phóng ra khi đốt cháy $1800$ gam than đá: $393{,}5.135 + 296{,}8.0{,}675 = 53282{,}34kJ$ \par Nhiệt lượng trên tương đương với số điện: $\dfrac {53282{,}34}{3600} \approx 14{,}8$ số.
\end{loigiaibth}
